% ==============================================================================
% T0 Theory / FFGFT: Document 274 (Chapter Content, DE)
% Titel: FFGFT (T0) im IPI-Feld -- Substrat- und Primitiv-Achsen
% Autor: Johann Pascher
% Datum: Juni 2026
% Referenzen (FFGFT-Vergleichsreihe): Dok. 244 (UMF), 245/269 (RA/PMT),
%   246 (RSG), 251 (Vopson), 260 (UC), 271 (HLV), 272 (Dot-Governance),
%   270 (Projektionskette), 190 P35 (xi^N-Trivialitaet).
% ==============================================================================

\section*{Abstract}
Diese Arbeit ordnet die FFGFT (T0, Zeit-Masse-Dualität) in das Feld der im
\emph{Information Physics Institute} (IPI) diskutierten Rahmenwerke ein. Statt eines
weiteren paarweisen Vergleichs (die liegen in Dok.~244, 245/269, 246, 251, 260, 271, 272 vor)
wird hier eine \emph{Landkarte} entlang zweier orthogonaler Ordnungsachsen vorgeschlagen:
(1)~die \textbf{Substrat-Geometrie} -- woraus der diskrete Untergrund besteht -- und
(2)~das \textbf{Primitiv} -- was als das Fundamentale gilt. Auf Achse~1 nimmt FFGFT den
\emph{periodischen} Pol ein (kompakter $T^4$-Kristall), HLV den \emph{aperiodischen}
(Quasikristall); der \emph{zufällige} Pol bleibt ein Kontrollmodell ohne festen Vertreter.
Diese Triade ist genau die, die der spektrale Diskriminator-Testbed (Dok.~271, beiliegendes
Skript) messbar macht. Auf Achse~2 steht FFGFT geometrie-primitiv neben informations-primitiven
(Vopson; stromabwärts Leavitt), algebraisch-hilbertraumbasierten (UMF) und selektions-/
constraint-orientierten Rahmenwerken (RA/PMT, RSG, UC, UIFT). Es wird ausdrücklich betont:
diese Zwei-Achsen-Karte ist eine \emph{Ordnungs-Überlagerung des Autors} -- in Dot-Theory-Sprache
ein \emph{sekundäres Rendering (O2)}, eine erklärte Lesart -- nicht das Eigenvokabular der
jeweiligen Rahmenwerke; das primäre Rendering (O0) verbleibt bei deren Urhebern.

% ============================================================================
\section{Zweck und Grenzen dieser Karte}
% ============================================================================

Im IPI-Austausch sind FFGFT in den letzten Monaten eine Reihe von Rahmenwerken
gegenübergestellt worden. Jede Gegenüberstellung endete auffallend oft mit demselben
Befund: \emph{Kategorieunterschied, nicht Widerspruch}. Dieses Dokument fragt, warum,
und schlägt eine Karte vor, die das erklärt.

Drei Vorbehalte vorab. \emph{Erstens}: die Etiketten unten (\enquote{periodisches Kristallgitter},
\enquote{informations-primitiv} usw.) sind \emph{Ordnungsbegriffe dieser Arbeit}. Manche
Urheber würden ihr Rahmenwerk anders einsortieren; die Karte ist eine Lesart, kein Urteil.
\emph{Zweitens} -- und das ist der wichtigere Punkt: die \emph{Wahl der beiden Achsen}
(Substrat-Geometrie, Primitiv) ist selbst eine \emph{deklarierte Projektion}, nicht eine
universelle Dimension des Feldes. Eine andere Untersucherin könnte ebenso zulässig andere
Achsen wählen -- etwa Zugänglichkeit/Constraint oder Beobachterabhängigkeit/Rückgewinnbarkeit
-- und käme zu einer anderen, gleichermaßen legitimen Klassifikation desselben Feldes. Die
hier gewählten Achsen werden daher als \emph{eigens ausgewiesene Projektion} geführt, nicht
als implizite Universalien; sie sind privilegiert allein relativ zum erklärten Zweck dieser
Karte. \emph{Drittens}: auch innerhalb dieser Projektion sortieren die Achsen das Feld nicht
restlos -- Achse~1 (Substrat) trennt sauber nur die geometrischen Substrat-Programme;
algebraische (UMF) und kontinuums-topologische ($\alpha$LGQV) Programme liegen neben der
Triade. Das ist kein Mangel der Karte, sondern ihre ehrliche Auflösung.

\medskip\noindent\textbf{Verwendung von O0/O2 (deklariert):} O0/O2 werden hier im
\emph{provenienzwahrenden} Sinn gebraucht -- als Marker dafür, wo die Deutungshoheit liegt
und wer das Korrekturrecht behält --, nicht als vollständiger Import von Guevaras
Rendering-Hierarchie (O0/O1/O2 als Repräsentationsebenen). Ob beide letztlich zusammenfallen,
bleibt offen; diese Karte verwendet allein die provenienzwahrende Lesart (vgl.\ Vossens
FAH-Anliegen). Rendering-Ebene und Provenienz fallen in diesem Dokument nur deshalb zusammen,
weil ein sekundäres Rendering \emph{fremder} Rahmenwerke notwendig von jemand anderem als dem
Urheber verfasst ist -- außerhalb des Fremdvergleichs entkoppeln sie.

% ============================================================================
\section{Achse 1 -- Substrat-Geometrie (die Kristall-Achse)}
% ============================================================================

FFGFT ist im Kern ein \emph{periodisches Kristallgitter}: der kompakte 4-Torus $T^4$ ist
periodisch und geordnet, mit fraktaler Korrektur $D_f=3-\xi$. HLV ist demgegenüber ein
\emph{aperiodischer Quasikristall} (Cut-and-Project $\mathbb{Z}^5\!\to\!3$D, golden ratio /
Fibonacci). Ein drittes, \emph{zufälliges} Substrat dient als Nullmodell. Diese drei Typen
sind nicht nur begrifflich verschieden -- sie haben \emph{verschiedene Spektral-Fingerabdrücke},
und genau diese Trennung macht der Graph-Laplace-Testbed aus Dok.~271 skaleninvariant messbar.

\begin{center}
\small
\begin{tabular}{@{}p{0.20\linewidth}p{0.22\linewidth}p{0.50\linewidth}@{}}
\hline
\textbf{Substrat-Typ} & \textbf{Vertreter} & \textbf{Spektraler Fingerabdruck (skaleninvariant)}\\
\hline
periodisches Kristallgitter
& \textbf{FFGFT} ($T^4$)
& stark entartet; $\langle r\rangle\!\approx\!0$; ganzzahlige $|n|^2$-Leiter (Summen von Quadraten); in $D{=}4$ wird \emph{jede} ganze Zahl getroffen (Lagrange); keine Cantor-Lücken\\
aperiodischer Quasikristall
& \textbf{HLV} (Krüger)
& geringe Entartung; \emph{intermediäres} (\enquote{kritisches}) $\langle r\rangle$; hohe Lücken-Reichhaltigkeit (Cantor-Hierarchie) -- erwartet, von beiden Kontrollen unterscheidbar\\
zufälliges Substrat
& Kontrollmodell
& $\langle r\rangle\!\approx\!0{,}50$ (GOE $0{,}531$); keine Entartung, keine Lückenhierarchie\\
\hline
\end{tabular}
\end{center}

\noindent Die Diskriminatoren ($\langle r\rangle$ = Verhältnis aufeinanderfolgender
Niveau-Abstände, Entartungsanteil, Lücken-Reichhaltigkeit) sind dimensionslos und brauchen
kein \emph{unfolding}; die absolute Skala spielt keine Rolle. Damit ist Achse~1 nicht bloß
ein Bild, sondern ein \emph{Experiment}: der Test, ob ein gegebenes Spektrum periodisch,
aperiodisch oder zufällig ist. FFGFT trägt dazu die Vorhersage des periodischen Arms bei
(die rationale $|n|^2$-Leiter), HLV das aperiodische Substrat; die Leitplanke
\enquote{$X=\varphi^N$ bzw.\ $X=\xi^N$ ist keine Vorhersage} (Dok.~190 P35) gilt für beide.

% ============================================================================
\section{Achse 2 -- das Primitiv (woraus alles folgt)}
% ============================================================================

Ein großer Teil des IPI-Feldes hat \emph{kein} geometrisches Substrat; dort greift die
Kristall-Achse nicht, und ein zweiter Ordnungsbegriff wird nötig: was gilt als das
Fundamentale?

\begin{itemize}
  \item \textbf{Geometrie-primitiv:} \textbf{FFGFT} -- alles folgt aus $T^4$-Geometrie und
        der einen dimensionslosen Konstante $\xi\approx 1/7500$. Statisch, beobachterunabhängig,
        ohne Selektor.
  \item \textbf{Informations-primitiv:} \textbf{Vopson} -- Masse-Energie-Informations-Äquivalenz
        (MEI), Zweites Gesetz der Infodynamik, Simulations-Hypothese. Information ist physikalisch
        und fundamental. Stromabwärts: \textbf{Leavitt}, dessen Bit-Budget-Konstruktion direkt auf
        Vopsons MEI aufsetzt.
  \item \textbf{Algebraisch / hilbertraumbasiert:} \textbf{UMF} (Gericke) -- ein primzahl-indizierter
        Hilbertraum $\ell^2(P_N)$ mit Kopplungs-Bivektor $B_{OI}$.
  \item \textbf{Selektions- / Constraint-orientiert} (was \emph{wählt oder beschränkt} die
        Konfiguration): \textbf{RA/PMT} (Guevara), \textbf{RSG} (Austin), \textbf{UC} (Haskins),
        \textbf{UIFT} (Teker).
  \item \textbf{Kontinuums-topologisch:} $\alpha$LGQV (Kriger) -- $S^3$-Topologie,
        Membran-Elastizität als Metrik.
  \item \textbf{Meta- / Governance} (ordnet \emph{andere} Theorien, eine Ebene höher):
        \textbf{Dot Theory} (Vossen).
\end{itemize}

Auf dieser Achse ist der FFGFT-Kontrast zu Vopson nicht periodisch-vs-aperiodisch, sondern
die \emph{geometrische gegen die informationelle Seite derselben MEI-Beziehung}: FFGFT nähert
sich von der Geometrie, Vopson von der Entropie (Dok.~251). Ein sichtbarer struktureller
Unterschied folgt daraus -- FFGFT schließt ein \enquote{Außen} bzw.\ eine Simulation
strukturell aus (es gibt keinen Außenstandpunkt zu einer statischen Gesamtstruktur), während
Vopsons Rahmen ihn offen lässt.

% ============================================================================
\section{Die Gesamtkarte}
% ============================================================================

\begin{center}
\footnotesize
\begin{tabular}{@{}p{0.155\linewidth}p{0.115\linewidth}p{0.205\linewidth}p{0.33\linewidth}p{0.07\linewidth}@{}}
\hline
\textbf{Rahmenwerk} & \textbf{Eigner} & \textbf{Kategorie (Achse)} & \textbf{Relation zu FFGFT} & \textbf{Dok.}\\
\hline
FFGFT / T0 & Pascher & Substrat: periodisch; Primitiv: Geometrie & --- (Bezugspunkt: statisch, beobachterfrei, kein Selektor) & ---\\
HLV & Krüger & Substrat: aperiodisch & gemeinsame spektral-/Rückgewinnungs-Achse; sonst Kategorienfehler & 271, 272\\
Vopson (MEI) & Vopson & Primitiv: Information & Geometrie- vs.\ Informations-Seite derselben MEI-Relation; Außen ausgeschlossen vs.\ offen & 251\\
(Bit-Budget) & Leavitt & Primitiv: Information (stromab Vopson) & teilt Vopsons Primitiv; Bit-Budget = $\Lambda$CDM-Werte als Eingabe, kein Vorwärts-Gehalt & ---\\
UMF & Gericke & Algebraisch / Hilbertraum & $\ell^2(P_N)$ ist echter Teilraum von $L^2(T^4)$; $\tau\!\cdot\!\mu=\hbar/c^2$ beidseitig & 244\\
RA/PMT & Guevara & Selektion (Meta-Schema) & FFGFT verweigert den Selektor-Slot RA1.5 (statisch) & 245, 269\\
RSG & Austin & Selektion / bulk-boundary & asymmetrische Komplementarität: FFGFT liefert Skala, RSG Selektionssprache & 246\\
UC & Haskins & Constraint (statisch) & beide statisch; $\alpha$-Korrektur-Inkonsistenz ($H_0$ vs.\ $a_0$) & 260\\
UIFT & Teker & holographischer Rand & Beobachter-Faktor-3 ist FFGFTs explizite Version des Volumen$\to$Rand-Moves & 267, 268\\
$\alpha$LGQV & Kriger & Kontinuums-topologisch & anderer Substrat-Typ ($S^3$, Membran); kein Gitter & ---\\
Dot Theory & Vossen & Meta / Governance & ordnet das Feld (Lexicon/OAP/Matrix/FAH); kein physikalisches Substrat & 272\\
\hline
\end{tabular}
\end{center}

\noindent\textbf{Abstammung und Außenkontakte} (keine IPI-Forum-Rahmenwerke, zur
Vollständigkeit): die \emph{Scale Relativity} (Nottale) ist der fraktale-Raumzeit-Vorfahr
der FFGFT; der Koide-Kreis (Koide, Rivero, Baez, Kocik, Brannen) betrifft die Massenformel-
Verifikation, nicht das Substrat.

% ============================================================================
\section{Warum so oft \enquote{Kategorieunterschied, nicht Widerspruch}}
% ============================================================================

Die Karte erklärt das wiederkehrende Vergleichsergebnis. FFGFT sitzt auf \emph{beiden}
Achsen am gleichen, strukturell-geometrischen Pol: periodischer Kristall (Achse~1) und
Geometrie-Primitiv, statisch, ohne Selektor (Achse~2). Die meisten anderen Rahmenwerke
bauen ihre Theorie gerade um einen \emph{Slot}, den FFGFT prinzipiell leer lässt oder anders
besetzt:

\begin{itemize}
  \item den \textbf{Selektor} (RA/PMT RA1.5, RSGs Selektionsprinzip) -- FFGFT ist statisch,
        es wird nichts ausgewählt;
  \item die \textbf{Volumen$\to$Rand-Naht} (UIFT) -- FFGFT trifft sie explizit als
        Beobachter-Faktor-3, löst aber nicht dieselbe offene Frage;
  \item den \textbf{Constraint} (UC) -- FFGFT instanziiert Struktur, formuliert aber kein
        eigenständiges Constraint-Prinzip;
  \item das \textbf{Informations-Primitiv} (Vopson) -- FFGFT nimmt die Geometrie als Primitiv
        und gewinnt Information daraus, nicht umgekehrt;
  \item den \textbf{Außenstandpunkt / die Simulation} (Vopson offen) -- in FFGFT strukturell
        ausgeschlossen.
\end{itemize}

Deshalb sind die Befunde selten Widersprüche: die Rahmenwerke reden über verschiedene Schichten.
Wo sie sich wirklich \emph{berühren} und damit empirisch kollidieren können, ist die
Substrat-Achse -- und dort liefert der spektrale Diskriminator (Dok.~271) den einzigen
derzeit scharfen, gemeinsamen Test.

% ============================================================================
\section{Anwendung der Governance-Schicht: vom Matrix-Eintrag zum A*}
% ============================================================================

In der Governance-Sprache von Vossen (Dot Theory) und Guevara (RA/PMT) lässt sich diese Karte
nicht nur als statische Einordnung lesen, sondern als \emph{Arbeitsablauf}. Die FFGFT--HLV-Achse
ist dafür das konkrete Beispiel -- der Ablauf ist bereits durchlaufen, nicht bloß beschrieben:

\begin{itemize}
  \item \textbf{Deklarierte Projektion:} die Substrat-Achse.
  \item \textbf{PRC} (Projection-Relative Classification): FFGFT periodisch, HLV aperiodisch.
  \item \textbf{PRM} (Projection-Relative Measurement): der spektrale Diskriminator
        (Niveau-Abstands-Verhältnis $\langle r\rangle$, Lückenstruktur, spektrale Dimension)
        -- eine Messung, die \emph{relativ zu dieser Projektion} definiert und ohne sie sinnlos ist.
  \item \textbf{Residuum} (deklariert, nicht absorbiert): eine unabhängige Reproduktion des
        HLV-Spektrums ($\mathbb{Z}^6$-Schnitt-Projektion, $k$-NN-Laplace) zeigt, dass der
        Quasikristall auf der Leitstatistik $\langle r\rangle$ \emph{noch nicht} vom
        Zufalls-Kontrollmodell trennbar ist; die aperiodische Signatur überlebt nur schwach,
        in Lückenstruktur und Tiefmoden-Sektor. Dieses Residuum wird ausgewiesen, nicht weggerechnet.
  \item \textbf{Operational Accord:} das gemeinsame Testprotokoll (periodisches Nullmodell,
        expliziter Rückgewinnungs-Limes, skaleninvariante Verhältnisse), gebunden an die beiden
        Beitragenden und ohne die Behauptung, die Rahmenwerke seien identisch.
\end{itemize}

\noindent Das Resultat -- der geteilte Spektral-Testbed -- ist ein \emph{Cooperative Admissible
Space (A*)} im Sinne der jüngsten Vossen--Guevara-Lokalisierung: eine deklarierte Region, in der
FFGFT und HLV für einen erklärten Zweck \emph{operativ} kooperieren, ohne ontologische Identität.
Guevaras Beobachtung gilt dabei: der Accord ist selbst ein Objekt mit eigenen
Admissibilitäts-Bedingungen (den drei vereinbarten Kriterien), eigener Residuumsstruktur (was der
abgeglichene Lauf zeigt) und einem Revisionsoperator (Neulauf, sobald HLV sein gewichtetes Spektrum
liefert). \emph{Operativ} statt bloß deklariert ist er, weil die Messung scheitern kann -- das
obige Residuum belegt, dass der Raum echte Arbeit leistet. O0 und O2 sind hier nicht das Ziel,
sondern die Instrumente, durch die der Accord konstruierbar und überprüfbar wird.

\medskip\noindent\emph{Auch dieser Abschnitt ist eine erklärte Anwendung fremden Vokabulars
(Vossen, Guevara) und steht unter Korrekturvorbehalt der Urheber.}

% ============================================================================
\section{Ehrliche Residuen (deklariert)}
% ============================================================================

\begin{itemize}
  \item Diese Karte ist eine \emph{erklärte Lesart} des Autors (Dot-Theory-O2). Das primäre
        Rendering (O0) jedes Rahmenwerks verbleibt bei seinem Urheber; die Einordnungen oben
        sind als Matrix-Einträge im Sinne von Dok.~272 zu lesen, nicht als Zuschreibungen.
  \item Die \emph{Wahl der Achsen selbst} ist eine deklarierte Projektion (Abschnitt~1), kein
        universelles Raster; andere Achsen (Zugänglichkeit/Constraint, Beobachterabhängigkeit/
        Rückgewinnbarkeit) sind gleichermaßen zulässig und ergäben eine andere legitime Karte.
  \item Die Substrat-Triade (periodisch/aperiodisch/zufällig) ordnet sauber nur FFGFT, HLV
        und das Kontrollmodell. UMF (algebraisch) und $\alpha$LGQV
        (kontinuums-topologisch) liegen neben der Triade; ihre Einordnung erfolgt allein über
        Achse~2 bzw.\ den Substrat-\emph{Typ}, nicht über die Spektral-Gabel.
  \item Die Platzierung von Leavitt als \enquote{stromabwärts Vopson} folgt aus der Herkunft
        seiner Bit-Budget-Konstruktion (MEI), nicht aus einer eigenständigen Substrat-Aussage.
  \item Sollte ein Urheber seinen Slot operationalisieren -- etwa Guevara RA1.5 als Gleichung,
        oder Vopson eine isolierte, vorhergesagte, dimensionslose Verhältnis-Vorhersage --,
        wäre die betreffende Zeile neu zu bewerten. Die Karte ist datiert und vorläufig
        (Juni 2026).
\end{itemize}

\noindent\textbf{Bezug:} Dok.~244 (UMF/Hilbertraum), 245/269 (RA/PMT), 246 (RSG),
251 (Vopson), 260 (UC/Haskins), 267/268 (UIFT-Naht), 270 (Projektionskette $T^4\!\to\!T^0$),
271/272 (HLV, paarweise und Dot-Governance), 190 P35 ($\xi^N$-Trivialität).
HLV: Krüger, Zenodo 10.5281/zenodo.20596861 und 10.5281/zenodo.20275888.
Vopson: AIP Advances 2022/2023/2024. RA/PMT: Guevara 2026.
